% ---
% filename : ["electrotechnique_triphase_20260426_117_couplage_resistances_triphase.tex"]
% id : "electrotechnique_triphase_20260426_117"
% title : "Couplage de résistances en régime triphasé"
% domain : "Électrotechnique"
% subdomain : "Triphasé"
% tags : ["triphasé", "étoile", "triangle", "puissance active", "résistance", "ts"]
% difficulty : 2
% time_solve : 12  # Calcul : 2 (difficulte) * (3 questions * 2)
% has_octave : true
% octave_files : ["electrotechnique_triphase_20260426_117_couplage_resistances_triphase.m"]
% author : "om"
% ---
\documentclass[a4paper,11pt,answers,addpoints]{exam}

% ---------------------------------------------------------
% LANGUE, ENCODAGE, FONTS
% ---------------------------------------------------------
\usepackage[utf8]{inputenc}

% ---------------------------------------------------------
% GÉOMÉTRIE ET MISE EN PAGE
% ---------------------------------------------------------
\usepackage[left=1.7cm,right=1.5cm,top=2cm,bottom=1cm]{geometry}
%\usepackage{a4wide}
\usepackage{microtype}      % Meilleure répartition du texte
\usepackage{float}          % Positionnement [H]
\usepackage{needspace}      % Saut conditionnel intelligent

% Éviter les veuves/orphelines (optionnel, à décommenter si besoin)
% \widowpenalty=10000
% \clubpenalty=10000

% Espacement vertical flexible
\raggedbottom
\setlength{\parskip}{0.5em plus 0.2em minus 0.1em}
\setlength{\parindent}{0pt}


% ---------------------------------------------------------
% MATHÉMATIQUES
% ---------------------------------------------------------
\usepackage{amsmath, amssymb, bm, esvect}

% Vecteurs
\newcommand{\V}{\overrightarrow}
\def\vect#1{\overrightarrow{\strut#1}}
\providecommand{\Degrees}[0]{\ensuremath{^\circ}}

% ---------------------------------------------------------
% UNITÉS ET GRANDEURS (SIunitx)
% ---------------------------------------------------------
\usepackage{siunitx}
\sisetup{output-decimal-marker={,}}
\DeclareSIUnit \var {var}
\DeclareSIUnit \VA {VA}
\DeclareSIUnit \kVA {kVA}
\DeclareSIUnit[quantity-product={}] \degree{\text{\textdegree}}

% ---------------------------------------------------------
% GRAPHIQUES : TikZ + CircuitikZ + PGFPlots
% ---------------------------------------------------------
\usepackage{tikz}
\usetikzlibrary{
	arrows.meta,
	intersections,
	decorations.pathreplacing,
	calligraphy,
	positioning
}

\usepackage[european,straightvoltages,EFvoltages]{circuitikz}

\usepackage{tkz-fct}
\usepackage{pgfplots}
\pgfplotsset{compat=1.12}

% Symbole étoile-triangle personnalisé
\newcommand{\wye}{%
	\mathbin{\tikz[x=1ex,y=1ex]{%
			\draw[line width=.1ex] (0,0)--(45:1)--++(-45:1) (45:1)--++(0,1);
	}}%
}


\usepackage{sankey}

\pgfdeclarelayer{background}
\pgfdeclarelayer{foreground}
\pgfdeclarelayer{sankeydebug}
\pgfsetlayers{background,main,foreground,sankeydebug}


% ---------------------------------------------------------
% TABLEAUX
% ---------------------------------------------------------
\usepackage{tabularx}
\usepackage{tabu}

% ---------------------------------------------------------
% HYPERLIENS
% ---------------------------------------------------------
\usepackage[colorlinks=true,
menucolor=red,
breaklinks=true,
urlcolor=blue,
linkcolor=blue]{hyperref}

% ---------------------------------------------------------
% COMMANDES PERSONNELLES
% ---------------------------------------------------------
\newcommand\nomauteur{bdminasmorgul@protonmail.com}
\newcommand\dateTe{29.04.2026}
\newcommand\brancheTe{TS}

% ---------------------------------------------------------
% STYLE DE L'EXERCICE
% ---------------------------------------------------------
\pagestyle{headandfoot}
\firstpageheadrule
\runningheadrule

\firstpageheader{\brancheTe\ (\nomauteur)}{Élève : }{(\dateTe) \thepage/\numpages}
\runningheader{\brancheTe\ (\nomauteur)}{Élève : }{(\dateTe) \thepage/\numpages}

% Compteur en bas de page
\newcommand{\totalpointtts}{%
	\begin{tikzpicture}[remember picture, overlay]
		\node[anchor=south west, inner sep=0pt, outer sep=0pt]
		at ([xshift=0.5cm,yshift=0.5cm]current page.south west)
		{\rotatebox{90}{Total des points : \numpoints}};
	\end{tikzpicture}
}

\firstpagefooter{\totalpointtts}{}{}
\runningfooter{\totalpointtts}{}{}

% Réglages exam
\printanswers
\addpoints
\pointsinmargin
\bracketedpoints

% ---------------------------------------------------------
% LISTINGS (code)
% ---------------------------------------------------------
\usepackage{xcolor}
\usepackage{listings}

\lstdefinestyle{octavestyle}{
	language=Matlab,
	basicstyle=\ttfamily\small,
	keywordstyle=\color{blue},
	commentstyle=\color{green!50!black},
	stringstyle=\color{red},
	stepnumber=1,
	frame=single,
	breaklines=true,
	breakatwhitespace=true,
	tabsize=4
}

% ---------------------------------------------------------
% DOCUMENT
% ---------------------------------------------------------
\begin{document}
\unframedsolutions
\pointsinmargin
\bracketedpoints

\section*{Préparation TS PQ CFC ELMO, IELE, PELE}

\begin{questions}
\question
Trois résistances chauffantes identiques de $R = \qty{80}{[\Omega]}$ chacune sont raccordées à un réseau triphasé européen standard $3 \times \qty{400}{[V]} / \qty{230}{[V]}$. Un technicien mesure une puissance active totale consommée par l'ensemble de $P_{mes} = \qty{1,984}{[kW]}$.

\begin{parts}
	\part[3] Calculer la puissance totale théorique $P_{\Delta}$ de l'installation si les résistances étaient couplées en triangle.
	\part[3] Calculer la puissance totale théorique $P_{Y}$ de l'installation si les résistances étaient couplées en étoile.
	\part[3] En comparant ces résultats avec la mesure du technicien, déterminer le type de couplage actuel. Justifier la réponse par le calcul.
\end{parts}
\vspace{0.5em}
\needspace{55\baselineskip}
\begin{solution}
	\begin{align*}
		a) \quad &\text{En couplage triangle, chaque résistance est soumise à la tension entre phases } U = \qty{400}{[V]}. \\[6pt]
		P_{\Delta} &= 3 \cdot P_{1phase} = 3 \cdot \dfrac{U^2}{R} \\[6pt]
		P_{\Delta} &= 3 \cdot \dfrac{400^2}{80} = 3 \cdot \dfrac{160000}{80} = \qty{6000}{[W]} = \qty{6}{[kW]} \\[6pt]
		b) \quad &\text{En couplage étoile, chaque résistance est soumise à la tension de phase } U_{ph} = \qty{230}{[V]}. \\[6pt]
		P_{Y} &= 3 \cdot P_{1phase} = 3 \cdot \dfrac{U_{ph}^2}{R} \\[6pt]
		P_{Y} &= 3 \cdot \dfrac{230^2}{80} = 3 \cdot \dfrac{52900}{80} = \qty{1983,75}{[W]} = \qty{1,984}{[kW]} \\[6pt]
		c) \quad &\text{La puissance mesurée } P_{mes} = \qty{1,984}{[kW]} \text{ correspond exactement au calcul du couplage étoile.} \\[6pt]
		&\text{Conclusion : Les résistances sont raccordées en \textbf{étoile}.} \\
	\end{align*}
	
	\begin{verbatim}
		% Télécharger OCTAVE depuis https://wiki.octave.org/Using_Octave et l'installer
		% Puis copier-coller le code dans OCTAVE
		R = 80;        % [Ohm] Valeur d'une résistance
		U = 400;       % [V] Tension composée (entre phases)
		Uph = 230;     % [V] Tension simple (entre phase et neutre)
		P_mes = 1.984; % [kW] Puissance mesurée par le technicien
		% a) Puissance en couplage triangle
		P_delta_W = 3 * (U^2 / R);
		P_delta_kW = P_delta_W / 1000;
		printf("a) Puissance en triangle : P_delta = %.2f [kW]\n", P_delta_kW);
		% b) Puissance en couplage étoile
		P_star_W = 3 * (Uph^2 / R);
		P_star_kW = P_star_W / 1000;
		printf("b) Puissance en etoile : P_star = %.3f [kW]\n", P_star_kW);
		
		% c) Vérification du couplage
		if abs(P_star_kW - P_mes) < 0.01
		printf("c) Le couplage detecte est : ETOILE\n");
		elseif abs(P_delta_kW - P_mes) < 0.01
		printf("c) Le couplage detecte est : TRIANGLE\n");
		else
		printf("c) Mesure non conforme aux couplages standards.\n");
		end
	\end{verbatim}
	\end{solution}
\end{questions}
\end{document}